% !TeX root = main.tex
% !TeX program = xelatex
% Translation of the current Chinese workflow and analysis objectives.

\section{Authorization Workflow and Analysis Objectives}
\label{sec:authorization-workflow-objectives}

\subsection{Authorization Workflow}
\label{subsec:authorization-workflow}

Fig.~\ref{fig:saga-authorization-workflow} illustrates SAGA's authorization workflow. Let $P$ denote the Provider, $A$ the receiver, and $B$ the sender. Both agents are registered, with their identities bound to their credentials.

\begin{figure}[tb]
  \centering
  % !TeX root = saga_authorization_compact.tex
% Normal successful path, at native publication size: 80 mm x 95 mm.
% Symbols: k is A's OTK public key, t the encrypted ACT, K a local shared key.
% Source: SAGA v2 IV-D2/IV-E and saga/{agent.py,provider/provider.py}.
% K is never transmitted. No atomicity or original-retriever binding is implied.
\begin{tikzpicture}[
  x=1mm,y=-1mm,
  font=\fontsize{7.5}{9}\selectfont,
  actor/.style={draw=black,line width=0.4pt,fill=white,
    minimum width=18mm,minimum height=4.6mm,inner sep=0pt,
    font=\bfseries\fontsize{8}{9.2}\selectfont},
  life/.style={draw=black,densely dashed,line width=0.4pt},
  msg/.style={-{Latex[length=1.4mm,width=1mm]},line width=0.5pt},
  wire/.style={fill=white,inner xsep=0.4mm,inner ysep=0.15mm},
  local/.style={draw=black,line width=0.4pt,fill=white,
    align=center,inner xsep=0.6mm,inner ysep=0.55mm,anchor=north},
  senderstep/.style={fill=white,align=center,inner sep=0.3mm},
  phase/.style={anchor=west,fill=white,inner xsep=0.5mm,
    inner ysep=0.2mm,font=\bfseries\fontsize{8}{9.2}\selectfont}
]
  \path[use as bounding box] (0,0) rectangle (80,95);
  \draw[life] (10,4.6) -- (10,38.5);
  \draw[life] (40,4.6) -- (40,94.8);
  \draw[life] (70,4.6) -- (70,94.8);
  \node[actor] at (10,2.3) {Provider $P$};
  \node[actor] at (40,2.3) {Sender $B$};
  \node[actor] at (70,2.3) {Receiver $A$};
  \node[fill=white,inner sep=0.2mm] at (40,8)
    {Registered agents; $P$: signed OTKs $k$; $A$: private keys.};

  \draw[line width=0.4pt] (0,12.5) -- (80,12.5);
  \node[phase] at (0.5,12.5) {I. Policy evaluation and OTK issuance};
  \draw[msg] (40,18.6) -- node[wire,above=0.3mm]
    {OTK request for $A$} (10,18.6);
  \node[local,text width=34mm] at (18,20.3) {
    Check identity, ACP,\\
    retrieval quota, available $k$.\\
    Issue $k$; debit quota; remove $k$.
  };
  \draw[msg] (10,33.5) -- node[wire,above=0.3mm]
    {Signed $k$ + $A$ metadata} (40,33.5);
  \node[senderstep] at (40,36.6) {Verify retrieved material};

  \draw[line width=0.4pt] (0,41.2) -- (80,41.2);
  \node[phase] at (0.5,41.2) {II. OTK consumption and ACT issuance};
  \draw[msg] (40,47) -- node[wire,above=0.3mm]
    {Identity material, $k$} (70,47);
  \node[local,text width=32.6mm] at (62.5,48.8) {
    Check identity, local ACP;\\
    consume unused $k$; derive $K$.\\
    Bind ACT to $B$'s PAC;\\
    set expiry/quota; issue/store.
  };
  \draw[msg] (70,65) -- node[wire,above=0.3mm]
    {Encrypted ACT $t$} (40,65);
  \node[senderstep] at (40,68.2) {Derive $K$; decrypt $t$; retain $t$};

  \draw[line width=0.4pt] (0,73) -- (80,73);
  \node[phase] at (0.5,73) {III. ACT request processing};
  \draw[msg] (40,78.8) -- node[wire,above=0.3mm]
    {Request $(t,\mathrm{msg})$} (70,78.8);
  \node[local,text width=32.6mm] at (62.5,80.5) {
    Validate ACT: PAC,\\
    expiry and request quota.\\
    Debit quota; process request.
  };
  \draw[msg] (70,94) -- node[wire,above=0.3mm]
    {Response} (40,94);
  \node[anchor=north west,align=left,fill=white,
    text width=32mm,inner sep=0.4mm] at (0.5,82) {
    Repeat requests while\\
    the ACT remains valid.
  };
\end{tikzpicture}

  \caption{Authorization workflow of SAGA.}
  \label{fig:saga-authorization-workflow}
\end{figure}

The receiver's ACP (Agent Contact Policy) defines contact permission through AID (agent identifier) patterns and budgets. When multiple rules match, the most specific matching rule determines the budget. Agent $A$ uploads pre-generated, signed OTK public keys to $P$ and retains the corresponding private keys. The Provider verifies the sender's registered identity and TLS (Transport Layer Security) client certificate, determines its OTK retrieval quota from the contact policy, and maintains the remaining quota for each sender~\cite{syros2025sagaV2,sagaReferenceImpl2026}.

After verifying the retrieved material, $B$ establishes a TLS connection with $A$ and uses the OTK to request an ACT. Agent $A$ consumes the unused OTK, derives a shared key from the corresponding private key and $B$'s PAC (public access control key), and issues an ACT bound to that PAC with a validity period and an initial request quota. Agent $B$ derives the same key using its private access control key and the OTK, decrypts the ACT, and retains the ciphertext and token metadata~\cite{syros2025sagaV2,sagaReferenceImpl2026}.

Subsequent requests carry the ACT ciphertext and are subject to identity binding, expiration, and remaining request quotas. An ACT can be reused across connections, with identity, local contact permission, and ACT validity checked at each new connection. The sender and receiver maintain separate local request quotas, and connections using the same ACT share the receiver's remaining quota. The ACT request quota is independent of the OTK retrieval quota maintained by $P$~\cite{sagaReferenceImpl2026}.

\subsection{Analysis Objectives}
\label{subsec:analysis-objectives-assumptions}

We examine three aspects of this authorization workflow.

\textbf{ACP:} Contact policies determine a sender's communication permission and OTK retrieval quota from its identity and policy rules. We examine whether matching respects rule priorities and authorization decisions conform to the specified policy~\cite{syros2025sagaV2,sagaReferenceImpl2026}.

\textbf{OTK:} The Provider issues an OTK, which the receiver consumes when issuing an ACT. We examine the identity and authorization conditions governing OTK issuance, use, and ACT issuance, including issuance limits and one-time consumption.

\textbf{ACT:} An ACT restricts subsequent requests through identity binding, a validity period, and a request quota. We examine token validation, request processing, and quota updates to assess quota enforcement.
